\begin{recipe}{Pommes Fondante}
\label{recipe:pommes-fondante}
\kicker[Bijgerecht,Frans,Kerst]{Bijgerecht · Frans · kerst}
\index[register]{Pommes fondante}
\index[register]{Aardappelen}
\index[register]{Boter}

\meta{45 minuten \dvd Oven 210\,°C}

\marginimage{images/pommes-fondante}
\lettrine{P}{ommes} fondante zijn aardappelen die rechtop in de pan worden gebakken tot goudbruin, waarna ze zachtjes garen in bouillon. Romig vanbinnen, krokant vanbuiten. Een elegant bijgerecht voor vlees of vis.

\blockrule{Ingrediënten}
\begin{ingredients}
  \ing{10}{vastkokende aardappelen (gelijk formaat)}
  \ing{50 g}{ongezouten roomboter}
  \ing{3--4 el}{olijfolie}
  \ing{250--350 ml}{kippenbouillon}
  \ing{4--5 takjes}{verse tijm}
  \ing{2--3 teentjes}{knoflook (optioneel, geplet)}
  \ingb{zout en peper}
\end{ingredients}

\blockrule{Bereiding}
\tip{Geen ovenbestendige pan? Bak de aardappelen eerst in een gewone pan, doe ze daarna over in een ovenschaal.}
\begin{steps}
  \item Verwarm de oven voor op 210\,°C.
  \item Schil de aardappelen. Snijd ze in dikke cilinders van ca. 4--5 cm hoog en snijd de ronde kanten eraf zodat ze stabiel staan. Spoel ze af onder koud water om zetmeel te verwijderen en dep droog. Bestrooi met zout en peper.
  \item Verhit de olie in een ovenbestendige pan op middelhoog vuur. Voeg de boter toe. Bak de aardappelen rechtop staand aan de onderkant goudbruin (ca. 5--7 minuten). Draai ze om en bak de andere kant ook goudbruin.
  \item Voeg de tijm, knoflook en bouillon toe. De bouillon moet ongeveer twee derde van de hoogte van de aardappelen bereiken. Breng kort aan de kook.
  \item Zet de pan in de voorverwarmde oven en laat 25--35 minuten garen tot de aardappelen zacht zijn en het vocht grotendeels is opgenomen. Giet eventueel halverwege wat bouillon bij als het te droog wordt.
  \item Haal uit de oven, bestrooi eventueel met wat extra zeezout en serveer direct.
\end{steps}
\end{recipe}
